% !TeX root = ../papers/nc-p3.tex
\section{有限点対応を保つ加法群と四つの推移変換}
\label{sec:finite-additive-windows}
標数零の代数閉体上、四辺の点対応が有限なFrobenius窓を取る。
有限完全交叉の射影は閉埋込みで、二つの隣接対応は共有位置で同じ長さ20の像を持つ
（定理\ref{thm:finite-frobenius-low-opposite-as}の証明）。
各辺の関係六平面に階数1のテンソルはなく、どの端点での評価階数も3以上である。
各位置の三次式の制限は全切断への同型である（定理\ref{thm:finite-square-chart}の完全交叉計算）。
ここでは窓の射影安定群内の非自明な加法群を考え、四つの文字空間上の冪零生成元を $N_i$ とする。

\begin{proposition}\label{prop:finite-additive-jordan}
各 $N_i$ のJordan型は $(4),(3,1),(2,1,1)$ のいずれかである。
各位置の最大ブロックの末端を検出する点、すなわち
$[p_i]=\mathbb P(\operatorname{im}(N_i^*)^{a_i-1})$（$a_i$ は最大ブロック長）は
四周期点列を与える。型 $(4)$ と型 $(2,1,1)$ は一窓に共存しない。
\end{proposition}
\begin{proof}
一位置が自明ならグラフ性により隣の作用も点スキーム上恒等である。
射影変換 $g$ の固定方程式 $v\wedge gv$ は二次式なのでその位置の全射影空間でも零になる。
従って四位置すべてが自明となる。隣接文字を $A,B$、最大長を $a,b$ とし、
$D=N_A\otimes1+1\otimes N_B$、辺の関係を $R$ とする。
$D^{a+b-2}$ は $N_A^{a-1}\otimes N_B^{b-1}$ の非零倍である。
$A$ の最大ブロックが一個ならこの作用の各像は純関係なので零である。
従って $R(p_A)\subset\ker N_B^{b-1}$、評価階数から $B$ の最大ブロックも一個で等号である。
これが $p_B$ を一意な後続点とする。最大ブロックが二個なら全四位置が $(2,2)$ である。
この場合 $D^2 R$ は二文字の像二平面の積に入り、純テンソルを含まないため次元は高々1である
（射影直線は二次曲面と交わる）。両商 $A/\operatorname{im}N_A,B/\operatorname{im}N_B$ への
関係の射影は両 $N$ による同型を通じて同じ次元を持つ。二つの固定直線の積に高々一式しかなく有限性に反する。

型 $(4)$ があれば全支持はその固定点 $p$ 一個にある。その位置の射影余接空間で誘導される作用は
長さ3の一ブロックである。型 $(2,1,1)$ の正準点の余接空間では作用は零である。
同じ局所環の余接空間は両方の商なので両型があれば次元は1であり、環は $C=k[t]/(t^{20})$ となる。
実は型 $(4)$ の単独支持でこの環は不可能である。$p(z)=1$ となる長い鎖の始文字を取り
$\phi=Nz/z$、$C$ に誘導する微分を $f(t)\partial_t$ と書く。
$\ker(N|_{\Sym^2 V})$ の次元は2である（四次元既約表現の対称平方は次元7と3の和）。
どの不変二次式 $q$ も $z^2$ の係数が零である。$N^6$ でその係数だけを検出できるからである。
非零な $Q=q/z^2$ は従って消失位数 $s\ge1$ を持ち、$\phi Q\ne0$ は三次式の単射性による。
微分方程式 $fQ'=-2\phi Q$ の右辺が零でないので先頭項を比較できる。
$e=\operatorname{ord}\phi$、$d$ をその先頭係数とすると $\operatorname{ord} f=e+1$ で、
$f$ の先頭係数 $c$ は $sc=-2d$ を満たす。すべての非零不変二次式に同じ $s$ を強制するが、
二つの独立な切断の先頭係数を一次結合で消せるため矛盾する。
\end{proof}

\begin{theorem}\label{thm:finite-four-transvections}
四つの $N_i$ がすべて型 $(2,1,1)$ なら、有限点対応を持つFrobenius窓は
標準quadratic Koszul左右AS正則である。この窓は半安定で、唯一の多重安定極限は
定理\ref{thm:one-standard-sl2-classification}の交代型または対称型の $P=0$ である。
加法群の固定平面どうしで切った各辺は長さ6で、元の辺の全ての閉点を台に含む。
\end{theorem}
\begin{proof}
各文字空間を $kz\oplus U\oplus kx$、$Nz=x$、$\dim U=2$ と書く。
重みを順に $-1,0,1$ とし、$z,x$ を標準表現にして四位置同時の $\mathrm{SL}_2$ に拡張する。
辺のテンソル積で重み $h$ 以上を $F^h$ と書く。$D^2R=0$ から $zz$ の係数は常に零である。
$R\cap F^1$ の重み1の成分を $xU$ と $Ux$ のそれぞれへ射影するといずれも単射である。
片側と $xx$ だけなら純テンソルになるからである。
一方 $DR\subset R\cap F^1$ の両射影は全射である。実際、$zU$ と $Uz$ の係数を読む写像は
正準点での評価の像 $U\oplus kx$ によりそれぞれ全射で、$D$ はこれらをそのまま重み1へ送る。
従って $R\cap F^1=DR$ は二次元、$R\cap F^0=\ker(D|_R)$ は四次元である。
後者では $D$ が重み2以上に入って消え、逆に重み $-1$ 上で $D$ は単射である。
残る重み0の二本の初期項は $UU\oplus k(xz-zx)$ にある。

重み $-1$ の二つの代表とその $D$ 像、重み0の不変な二代表を基底にすると、
抽象的な $R$ 自身に標準表現二個と自明表現二個の構造を与えられる。
この構造での辺の埋込みは重みを下げず、次数零の部分も埋込みである。
反対辺間の窓の完全対合 $R_i^*\simeq R_{i+2}$ と逆写像は加法群同変である。
二つの $\mathrm{SL}_2$ 表現の間の加法群同変写像は $\operatorname{Hom}$ の最高重みベクトルの和なので
重みを下げず、次数零の部分は $\mathrm{SL}_2$ 同変である。逆写像との積により次数零の部分は同型である。
従って四位置の対角一径数群による窓の極限 $w_0$ は存在し、
両方の隣接切断で六次元の初期関係空間の間の完全対合を与える。
この極限は $\mathrm{SL}_2$ 不変で、標準型の重複度空間から両座標への射影は上記の二つの全射である。
したがって標準型切断の四つの $b_i$ は可逆である（同定理の射影による分解）。
各文字の支持も全体である。重み $-1$ と $1$ の混合関係の両射影から $U,z,x$ を全て検出できる。
同定理で $w_0$ は正則、命題\ref{prop:unrestricted-good-reduction}で元の窓も正則である。

二つの固定平面 $x=0$ 上で $DR$ の二式が消える。残り四式は
$\mathbb P^2\times\mathbb P^2$ 上で有限、かつ正準点対を含むので、独立な完全交叉で長さは
$\binom42=6$ である。連結群は有限支持の各点を固定する。
重み0の不変な二代表の制限はその $UU$ 成分のみで、これらも独立である。
$w_0$ が $P=0$ なら自明表現部分の切断の $UU$ 射影は $kb_i$ だけになり、矛盾する。
従って $w_0$ は交代型または対称型の $P\ne0$ であり、系\ref{cor:semisimple-symmetry-closed-limits}
から多重安定極限が従う。この必要条件は二つの型の有限例の存在を仮定しない。
\end{proof}
